\begin{proof}
By \Cref{cor:ch19-continuous-zero-current-harmonic}, the function \(h\) is
Thuillier harmonic.  Write \(\mathcal H_C\) for the sheaf of Thuillier
harmonic germs and \(\mathcal F_C\) for the real vector subsheaf generated
by germs \(\log|f|\) with \(f\) an invertible analytic function.  Thuillier's
description, recalled in \Cref{rem:ch08-harmonic-vs-log}, says that
\(\mathcal H_C/\mathcal F_C\) is supported at the type~\(2\) points and
that
\[
       (\mathcal H_C/\mathcal F_C)_x
       \simeq\operatorname{Pic}^0(C_x)\otimes_{\mathbf Z}\mathbf R
\]
there; see \cite[Theorem~2.3.21 and Lemma~2.3.22]{Thu05}.  Under this
isomorphism, the class of the germ \(h_x\) is precisely
\(\operatorname{ob}_x(h)\).

On the other hand, \cite[Remark~15.10]{GJR26} identifies
\(\mathcal F_C\) with the kernel of
\[
 \dprimeddprime:\mathcal A^{0,0}_{\mathrm{CLD}}
                    \longrightarrow\mathcal A^{1,1}_{\mathrm{CLD}},
\]
that is, with the sheaf of ordinary CLD-smooth harmonic functions.  Hence a
type~\(2\) harmonic germ lies in that ordinary smooth kernel exactly when
its residual Picard class vanishes.  Away from type~\(2\), the quotient
sheaf is zero, which proves the assertion for points of types \(1\), \(3\),
and \(4\) without imposing a nonexistent annular coordinate at a type~\(4\)
point.  Since ordinary smoothness and the smooth zero-Hessian equation are
local, \eqref{eq:p213-residual-picard-vanishing} proves the global assertion
on \(W\).

Finally, the two stated sufficient hypotheses for the quotient to vanish
are exactly the special cases of \cite[Theorem~2.3.21]{Thu05} recorded in
\Cref{rem:ch08-harmonic-vs-log}.  Equivalently, the abstract group of
degree-zero line-bundle classes on every \(C_x\) is torsion; no uniform
torsion exponent is required.
\end{proof}
